\section{Extraction de la Courbe des Taux et Nelson-Siegel}
\label{sec:taux}

\subsection{Parité Call-Put et Taux Implicite}

En l'absence d'arbitrage, la parité call-put donne pour un call $C$ et un put $P$
de même strike $K$ et maturité $T$ :
\begin{equation}
  C - P = e^{-rT}(F_T - K)
  \label{eq:put_call_parity}
\end{equation}
d'où le taux implicite continu :
\begin{equation}
  r(T) = -\frac{1}{T}\ln\!\left(\frac{C - P}{F_T - K}\right)
  \label{eq:taux_implicite}
\end{equation}

Pour robustesse, on calcule la médiane sur les paires $(C, P)$ valides
(même strike, même maturité) pour chaque $T$.

\subsection{Modèle Nelson-Siegel}

On lisse la courbe des taux par le modèle de Nelson-Siegel (1987) :
\begin{equation}
  r_{NS}(T) = \beta_0 + \beta_1\,\frac{1 - e^{-\lambda T}}{\lambda T}
            + \beta_2\!\left(\frac{1 - e^{-\lambda T}}{\lambda T} - e^{-\lambda T}\right)
  \label{eq:nelson_siegel}
\end{equation}

Les paramètres $(\beta_0, \beta_1, \beta_2, \lambda)$ sont calibrés par
minimisation des moindres carrés (évolution différentielle + multi-départ L-BFGS-B).

% TODO : insérer courbe taux empiriques vs NS (plot_rates_curve)

\subsection{Résultats}

% TODO : tableau des paramètres NS calibrés + RMSE, R²
